% Submission guidelines:
% The Title Page will not be sent to peer reviewers and must include:
% - Manuscript title
% - Authors’ names and affiliations
% - Address for correspondence including email addresses
% - Acknowledgments
% It must also be submitted in .doc or .docx format because professionalism means nothing to us.

\documentclass[10pt]{article}
\usepackage{hyperref}
\usepackage{xcolor}
\renewcommand{\thefootnote}{\arabic{footnote}}
\hypersetup{colorlinks = true, allcolors = {blue}}
\begin{document}
\pagenumbering{gobble}
\title{An FPGA-based Accelerator for 3D CNNs Using FFT convolution}
\author{
  Ricardo~Di~Curzio~Lera, 
  Miguel~Velasques~Abilio~Piola~Alves,\\
  and~Bruno~de~Carvalho~Albertini,~\textit{Fellow},~IEEE%
  \footnote{%
  R. Lera, M. Alves and B. Albertini are with the Polytechnic School of the University of São Paulo, São Paulo, Brazil.\protect\\
  E-mails: \href{mailto:ricardodclera@usp.br}{ricardodclera@usp.br}; \href{mailto:me@somewhere.com}{miguel.velasques@usp.br}; \href{mailto:me@somewhere.com}{balbertini@usp.br}\protect\\
  Research was funded by Itaú Unibanco S.A. and the Foundation of the University of São Paulo (FUSP); no limitations apply.%
}}
\date{}
\maketitle
\begin{abstract}
  Accelerating 3D CNNs poses a significant design challenge, due in large part to the increase in arithmetic complexity of the convolution operation in three dimensions. Winograd-based designs can achieve higher resource efficiency in networks using small kernels, yet are limited by increasingly complex transform operations in higher dimensions. FFT convolution has seen recent success in applying overlap-add to large CNNs. The Complex-domain multiplications required by the algorithm can likewise be optimized using polar coordinate representation. We propose a 3D CNN accelerator using FFT convolution augmented by overlap-add and CORDIC algorithms. We implement and test this accelerator on FPGA technology, achieving over double DSP-efficiency compared to the state-of-the-art.
\end{abstract}
\textbf{Index Terms} ---
  Convolutional neural networks,
  % Neural network hardware,
  % Open source hardware,
  Fast Fourier transforms,
  Computational efficiency,
  Algorithmic efficiency,
  Field programmable gate arrays,
  % AI accelerators,
  Accelerator architectures,
  Hardware acceleration.
\end{document}